%=============================================================
%  DATOS DE LA BÚSQUEDA Y DEL CRIBADO
%  Las fechas, los conteos por base, el flujo PRISMA y los kappa los
%  escribe automáticamente cribado/scripts/06_resultados_prisma.py.
%  A mano solo: \RegistroOSF y \RepositorioDatos.
%  Lo pendiente aparece en rojo en el PDF.
%=============================================================

% --- Fechas y registro ---------------------------------------
\newcommand{\FechaBusqueda}{19/10/2026}
\newcommand{\FechaBusquedaEN}{19 October 2026}
\newcommand{\RegistroOSF}{\pend{OSF, DOI 10.17605/OSF.IO/XXXXX}}
\newcommand{\RepositorioDatos}{\url{https://github.com/ClaudioB12/rsl-wilbert}}

% --- Revisor y fiabilidad del cribado (autor único) ------------
\newcommand{\Revisor}{W.\,A.\,M.\,A.}
\newcommand{\KappaRecribadoTA}{\pend{0,xx}}     % re-cribado ciego, título y resumen
\newcommand{\NRecribadoTA}{\pend{n}}
\newcommand{\KappaRecribadoTC}{\pend{0,xx}}     % re-cribado ciego, texto completo
\newcommand{\NRecribadoTC}{\pend{n}}
\newcommand{\KappaAsistenteTA}{0,06}     % acuerdo autor-asistente (TA)
\newcommand{\AcuerdoAsistenteTA}{31~\%}
\newcommand{\NExclRegla}{100}             % exclusiones propuestas por regla
\newcommand{\NExclReglaConfirmadas}{100}  % ... confirmadas por el autor
\newcommand{\KappaIATA}{0,77}            % acuerdo autor-IA en TA (pasa / no pasa)
\newcommand{\AcuerdoIATA}{93~\%}
\newcommand{\NCambiosIATA}{62}            % decisiones de la IA que el autor cambió
\newcommand{\NFueraFiltros}{123}           % quitados por filtros del protocolo aplicados por script
\newcommand{\NexclOA}{577}                 % excluidos en TA por no ser de acceso abierto (CI-2)
\newcommand{\FechaOA}{19/10/2026}             % fecha de verificación del acceso abierto
\newcommand{\AcuerdoCalidad}{\pend{xx~\%}}     % puntaje propuesto vs. verificado

% --- Registros por base (Tabla de bases) ----------------------
\newcommand{\NScopus}{498}
\newcommand{\NWoS}{193}
\newcommand{\NIEEE}{152}
\newcommand{\NScienceDirect}{337}
\newcommand{\NACM}{397}
\newcommand{\FScopus}{19/10/2026}
\newcommand{\FWoS}{19/10/2026}
\newcommand{\FIEEE}{19/10/2026}
\newcommand{\FScienceDirect}{19/10/2026}
\newcommand{\FACM}{19/10/2026}

% --- Flujo PRISMA 2020 ----------------------------------------
\newcommand{\Pidentificados}{1577}   % suma de las cinco bases
\newcommand{\Pduplicados}{337}
\newcommand{\Pcribados}{1240}
\newcommand{\PexclTA}{1137}
\newcommand{\Pbuscados}{103}        % solicitados a texto completo
\newcommand{\PnoRecuperados}{11}
\newcommand{\Pevaluados}{92}
\newcommand{\PexclTC}{20}
% Motivos de exclusión a texto completo (deben sumar \PexclTC)
\newcommand{\PexclCEdos}{2}      % CE-2 no es tráfico vehicular
\newcommand{\PexclCEtres}{7}     % CE-3 solo pronóstico/clasificación
\newcommand{\PexclCEcuatro}{2}   % CE-4 sin DL como método principal
\newcommand{\PexclCEcinco}{3}    % CE-5 tipo de documento / preprint
\newcommand{\PexclCIseis}{5}     % CI-6 sin evaluación cuantitativa de la detección
\newcommand{\PexclCEseis}{0}     % CE-6 información insuficiente
\newcommand{\PexclOtros}{1}      % CE-1, CI-2, CI-3
% Otros métodos (búsqueda v1 en ACM DL a texto completo)
\newcommand{\POtrosIdentificados}{0}
\newcommand{\POtrosIncluidos}{0}

% --- Base de recuperación de cada estudio (Anexo B) -----------
% La escribe cribado/scripts/06_resultados_prisma.py en generado/base_recuperacion.tex.
\newcommand{\BaseRecup}[1]{\ifcsname baserecup#1\endcsname\csname baserecup#1\endcsname\else\pend{base}\fi}
